\documentclass[11pt]{article}

\usepackage[margin=1in]{geometry}
\usepackage{amsmath}
\usepackage{amssymb}
\usepackage{graphicx}
\usepackage{array}
\usepackage{booktabs}
\usepackage{enumitem}
\usepackage{fancyvrb}
\usepackage[table]{xcolor}
\usepackage{microtype}
\usepackage{titlesec}
\usepackage{fancyhdr}
\usepackage{forest}

% --- Disjoint-set forest node style ---
\forestset{
  dsf/.style={
    for tree={draw, rounded corners, semithick, minimum width=2.4em,
      minimum height=1.9em, inner sep=2pt, font=\small, anchor=center,
      edge={semithick}, s sep=5mm, l sep=7mm,
      parent anchor=south, child anchor=north},
  },
}

% --- Colour palette ---
\definecolor{accent}{HTML}{1F4E79}   % deep blue for headings / boxes

% --- Section & subsection styling (with a rule under each section) ---
\titleformat{\section}
  {\normalfont\large\bfseries\color{accent}}{}{0pt}{}[{\color{accent!40}\titlerule[1pt]}]
\titlespacing*{\section}{0pt}{1.6em}{0.7em}
\titleformat{\subsection}
  {\normalfont\bfseries\color{accent!85!black}}{}{0pt}{}
\titlespacing*{\subsection}{0pt}{0.9em}{0.3em}

% --- Slightly looser table rows ---
\renewcommand{\arraystretch}{1.25}

% --- Highlighted final-answer box ---
\newcommand{\answer}[1]{%
  \par\smallskip\noindent
  \fcolorbox{accent}{accent!6}{%
    \parbox{\dimexpr\linewidth-2\fboxsep-2\fboxrule\relax}{%
      \textbf{\color{accent}Answer.}\ #1}}%
  \par\smallskip}

% --- Running header / footer ---
\pagestyle{fancy}
\fancyhf{}
\renewcommand{\headrulewidth}{0.4pt}
\lhead{\small COMP9312 --- Assignment 1}
\rhead{\small Solutions}
\cfoot{\small\thepage}

\title{\vspace{-1.2em}\textbf{COMP9312 --- Assignment 1 Solutions}}
\author{YUHUI LUO \quad (z5565446)}
\date{}

\begin{document}
\maketitle
\thispagestyle{fancy}

%=====================================================================
\section*{Q1. Disjoint-set forests for connected components}

\textbf{Setup.} Union by size: the larger tree's root stays as the root and the
smaller tree's root is attached under it. \textbf{Tie-break:} when both trees have
equal size, the tree whose root is \emph{lexicographically smaller} is treated as
larger, so its root becomes the new root. (No path compression is applied.)

Notation: each node is a set element; a root is labelled \texttt{root\,(size)} and its
children hang below it. Initially eight singletons $A,B,C,D,E,F,G,H$, each of size $1$.
Each drawing shows only the non-trivial trees; any vertex not shown remains an isolated
singleton.

\medskip
\noindent\begin{minipage}{\linewidth}
\noindent\textbf{Edge 1: B--C} --- both size $1$, tie $\Rightarrow$ smaller root $B$ wins.
\begin{center}
\begin{forest} dsf
[{$B$\,(2)} [{$C$}]]
\end{forest}
\end{center}
\noindent Components: $\{B,C\}$; singletons $A,D,E,F,G,H$.
\end{minipage}

\medskip
\noindent\begin{minipage}{\linewidth}
\noindent\textbf{Edge 2: A--C} --- $A$ (size $1$) vs.\ tree rooted $B$ (size $2$); $B$ larger.
\begin{center}
\begin{forest} dsf
[{$B$\,(3)} [{$C$}] [{$A$}]]
\end{forest}
\end{center}
\noindent Components: $\{A,B,C\}$; singletons $D,E,F,G,H$.
\end{minipage}

\medskip
\noindent\begin{minipage}{\linewidth}
\noindent\textbf{Edge 3: D--E} --- both size $1$, tie $\Rightarrow$ smaller root $D$ wins.
\begin{center}
\begin{minipage}[t]{0.46\textwidth}\centering
\begin{forest} dsf
[{$B$\,(3)} [{$C$}] [{$A$}]]
\end{forest}
\end{minipage}\hfill
\begin{minipage}[t]{0.46\textwidth}\centering
\begin{forest} dsf
[{$D$\,(2)} [{$E$}]]
\end{forest}
\end{minipage}
\end{center}
\noindent Components: $\{A,B,C\}$, $\{D,E\}$; singletons $F,G,H$.
\end{minipage}

\medskip
\noindent\begin{minipage}{\linewidth}
\noindent\textbf{Edge 4: C--E} --- tree $B$ (size $3$) vs.\ tree $D$ (size $2$); $B$ larger,
so $D$ attaches to $B$.
\begin{center}
\begin{forest} dsf
[{$B$\,(5)} [{$C$}] [{$A$}] [{$D$} [{$E$}]]]
\end{forest}
\end{center}
\noindent Components: $\{A,B,C,D,E\}$; singletons $F,G,H$.
\end{minipage}

\medskip
\noindent\begin{minipage}{\linewidth}
\noindent\textbf{Edge 5: F--G} --- both size $1$, tie $\Rightarrow$ smaller root $F$ wins.
\begin{center}
\begin{minipage}[t]{0.6\textwidth}\centering
\begin{forest} dsf
[{$B$\,(5)} [{$C$}] [{$A$}] [{$D$} [{$E$}]]]
\end{forest}
\end{minipage}\hfill
\begin{minipage}[t]{0.34\textwidth}\centering
\begin{forest} dsf
[{$F$\,(2)} [{$G$}]]
\end{forest}
\end{minipage}
\end{center}
\noindent Components: $\{A,B,C,D,E\}$, $\{F,G\}$; singleton $H$.
\end{minipage}

\medskip
\noindent\begin{minipage}{\linewidth}
\noindent\textbf{Edge 6: E--G} --- tree $B$ (size $5$) vs.\ tree $F$ (size $2$); $F$ attaches
to $B$.
\begin{center}
\begin{forest} dsf
[{$B$\,(7)} [{$C$}] [{$A$}] [{$D$} [{$E$}]] [{$F$} [{$G$}]]]
\end{forest}
\end{center}
\noindent Components: $\{A,B,C,D,E,F,G\}$; singleton $H$.
\end{minipage}

\medskip
\noindent\begin{minipage}{\linewidth}
\noindent\textbf{Edge 7: H--F} --- tree $B$ (size $7$) vs.\ $H$ (size $1$); $H$ attaches to $B$.
\begin{center}
\begin{forest} dsf
[{$B$\,(8)} [{$C$}] [{$A$}] [{$D$} [{$E$}]] [{$F$} [{$G$}]] [{$H$}]]
\end{forest}
\end{center}
\noindent Components: $\{A,B,C,D,E,F,G,H\}$ --- a single component.
\end{minipage}

\medskip
\noindent\begin{minipage}{\linewidth}
\noindent\textbf{Edge 8: G--H} --- $\mathrm{find}(G)=\mathrm{find}(H)=B$: already in the same
set $\Rightarrow$ \textbf{no change}; the forest is identical to that after Edge~7.
\begin{center}
\begin{forest} dsf
[{$B$\,(8)} [{$C$}] [{$A$}] [{$D$} [{$E$}]] [{$F$} [{$G$}]] [{$H$}]]
\end{forest}
\end{center}
\noindent Components: $\{A,B,C,D,E,F,G,H\}$ --- unchanged (still one component).
\end{minipage}

\answer{All eight elements end up in a single connected component: one tree of size~$8$ rooted at $B$, as shown after Edge~7.}

%=====================================================================
\section*{Q2. BFS on a grid map}

Start $S=(1,1)$. Neighbour scan order \textbf{Right $\to$ Down $\to$ Left $\to$ Up}.
Visited time $=$ the dequeue order, starting at $0$.

\subsection*{(a) Visited times}

Dequeue sequence (\texttt{time : cell}):
\begin{verbatim}
0:(1,1)  1:(1,2)  2:(1,3)  3:(2,3)  4:(3,3)  5:(3,4)  6:(3,2)
7:(3,5)  8:(3,1)  9:(3,6)  10:(4,5) 11:(4,1) 12:(5,5) 13:(5,1)
14:(5,6) 15:(6,5) 16:(5,2) 17:(6,6) 18:(6,4) 19:(5,3) 20:(6,3)
\end{verbatim}

Grid with visited times (\texttt{\#} $=$ blocked, \texttt{inf} $=$ passable but
unreachable from $S$):
\begin{verbatim}
       c1  c2  c3  c4  c5  c6
   r1:  0   1   2   #   inf inf
   r2:  #   #   3   #   #   #
   r3:  8   6   4   5   7   9
   r4:  11  #   #   #   10  #
   r5:  13  16  19  #   12  14
   r6:  #   #   20  18  15  17
\end{verbatim}

Cells $(1,5)$ and $(1,6)$ are passable but walled off from $S$, so BFS never reaches
them.

\subsection*{(b) Queries (arbitrary neighbour order)}

A FIFO BFS dequeues vertices in non-decreasing \textbf{distance from $S$} (all of
level $k$ before any of level $k+1$), regardless of neighbour order. Hence:
\begin{itemize}[nosep]
  \item $\mathrm{dist}(x)<\mathrm{dist}(y)\Rightarrow x$ always before $y\Rightarrow$ \textbf{true};
  \item $\mathrm{dist}(x)>\mathrm{dist}(y)\Rightarrow x$ always after $y\Rightarrow$ \textbf{false};
  \item $\mathrm{dist}(x)=\mathrm{dist}(y)\Rightarrow$ a neighbour order can put either first $\Rightarrow$ \textbf{false}.
\end{itemize}

Distances from $S$:
\begin{verbatim}
(1,1)=0 (1,2)=1 (1,3)=2 (2,3)=3 (3,3)=4
(3,2)=5 (3,4)=5
(3,1)=6 (3,5)=6
(4,1)=7 (3,6)=7 (4,5)=7
(5,1)=8 (5,5)=8
(5,2)=9 (5,6)=9 (6,5)=9
(5,3)=10 (6,6)=10 (6,4)=10
(6,3)=11
\end{verbatim}

\begin{center}
\rowcolors{2}{accent!6}{white}
\begin{tabular}{c c c c c l}
\toprule
\rowcolor{white}
Query & $x$ & $y$ & $\mathrm{dist}(x)$ & $\mathrm{dist}(y)$ & Result \\
\midrule
1 & $(1,2)$ & $(3,2)$ & 1  & 5  & \textbf{true} \\
2 & $(3,4)$ & $(3,2)$ & 5  & 5  & \textbf{false} (same level) \\
3 & $(5,1)$ & $(3,6)$ & 8  & 7  & \textbf{false} ($x$ deeper) \\
4 & $(4,5)$ & $(5,6)$ & 7  & 9  & \textbf{true} \\
5 & $(6,6)$ & $(5,3)$ & 10 & 10 & \textbf{false} (same level) \\
6 & $(3,1)$ & $(5,5)$ & 6  & 8  & \textbf{true} \\
7 & $(5,6)$ & $(4,5)$ & 9  & 7  & \textbf{false} ($x$ deeper) \\
8 & $(6,4)$ & $(6,3)$ & 10 & 11 & \textbf{true} \\
\bottomrule
\end{tabular}
\end{center}

%=====================================================================
\section*{Q3. Graph representation and algorithm time complexity}

Assume a vertex-indexed \texttt{visited}/\texttt{dist} array of size $n$ ($O(1)$
access, $O(n)$ to initialise).

Reasoning per operation:
\begin{itemize}[nosep]
  \item \textbf{SingleSourceShortestDistances$(u)$:} unweighted $\Rightarrow$ BFS; cost
    $=O(n)$ init $+$ enumerating the neighbours of every visited vertex.
  \item \textbf{CountConnectedComponents$()$:} a full BFS/DFS sweep over all vertices and edges.
  \item \textbf{CountCommonNeighbours$(u,v)$:} enumerate the neighbours of one vertex and
    test membership in the other (cost driven by the per-structure membership mechanism).
  \item \textbf{DFS$(u)$:} same enumeration cost as the BFS sweep from $u$.
\end{itemize}

The only structure that cannot enumerate neighbours in $O(\deg)$ is the adjacency
matrix (a row scan is always $O(n)$), so all of its traversals become $O(n^2)$.

For brevity write $d_u=\deg(u)$ and $d_v=\deg(v)$.

\medskip
\noindent\textbf{Traversal-based operations.} All three are a single graph sweep, so they
share the same cost --- only the per-vertex neighbour enumeration differs between structures:
\begin{center}
\rowcolors{2}{accent!6}{white}
\begin{tabular}{l c c c}
\toprule
\rowcolor{white}
 & \textbf{SSSD$(u)$} & \textbf{CountCC$()$} & \textbf{DFS$(u)$} \\
\midrule
1. Adjacency matrix & $O(n^2)$ & $O(n^2)$ & $O(n^2)$ \\
2. Adjacency arrays (unsorted) & $O(n+m)$ & $O(n+m)$ & $O(n+m)$ \\
3. Sorted adjacency arrays & $O(n+m)$ & $O(n+m)$ & $O(n+m)$ \\
4. Hash-set adjacency lists & $O(n+m)$ & $O(n+m)$ & $O(n+m)$ \\
5. CSR & $O(n+m)$ & $O(n+m)$ & $O(n+m)$ \\
6. Adjacency balanced BSTs & $O(n+m)$ & $O(n+m)$ & $O(n+m)$ \\
\bottomrule
\end{tabular}
\end{center}

\begin{minipage}{\linewidth}
\noindent\textbf{Common-neighbour query.} This is the operation that distinguishes the
structures, since it depends on each one's membership-test mechanism:
\begin{center}
\rowcolors{2}{accent!6}{white}
\begin{tabular}{l l}
\toprule
\rowcolor{white}
 & \textbf{CountCommonNeighbours$(u,v)$} \\
\midrule
1. Adjacency matrix & $O(n)$ \\
2. Adjacency arrays (unsorted) & $O(d_u d_v)$ \\
3. Sorted adjacency arrays & $O(d_u+d_v)$ \\
4. Hash-set adjacency lists & $O(\min(d_u,d_v))$ \\
5. CSR & $O(d_u d_v)$ \\
6. Adjacency balanced BSTs & $O(\min(d_u,d_v)\,\log\max(d_u,d_v))$ \\
\bottomrule
\end{tabular}
\end{center}
\end{minipage}

\noindent\textbf{Notes:}
\begin{itemize}[nosep]
  \item \textbf{Matrix CommonNeighbours $=O(n)$:} scan $w=1\ldots n$, test $M[u][w]\wedge M[v][w]$.
  \item \textbf{Unsorted arrays / CSR:} for each of $u$'s neighbours, a membership test in
    $v$'s list is a linear scan $O(d_v)\Rightarrow O(d_u d_v)$.
  \item \textbf{Sorted arrays:} merge the two sorted neighbour lists $\Rightarrow O(d_u+d_v)$.
  \item \textbf{Hash sets (traversals):} worst-case $O(n+m)$ --- BFS/DFS only \emph{iterate}
    each neighbour set ($O(\deg)$) and test the vertex-indexed \texttt{visited} array ($O(1)$);
    no hashing lookups are involved.
  \item \textbf{Hash sets (CommonNeighbours):} iterate the smaller set and do
    $O(1)$-expected lookups in the larger $\Rightarrow O(\min(d_u,d_v))$
    \emph{expected} (the only operation that relies on hashing).
  \item \textbf{Balanced BST (CommonNeighbours):} iterate the smaller neighbour set and search each
    element in the larger tree ($O(\log)$ per search) $\Rightarrow
    O(\min(d_u,d_v)\,\log\max(d_u,d_v))$.
\end{itemize}

%=====================================================================
\section*{Q4. Forced order in all possible DFS traversals}

\subsection*{(a) Pseudocode}

\textbf{Characterization.} $u$ is visited before $v$ in \emph{every} DFS from $s$
\textbf{iff every path from $s$ to $v$ passes through $u$} (i.e.\ $u$ \emph{dominates} $v$).
\begin{itemize}[nosep]
  \item \emph{Sufficient:} if every $s$--$v$ path uses $u$, the DFS-tree path from $s$ to
    $v$ also uses $u$, so $u$ (an ancestor of $v$) is discovered first.
  \item \emph{Necessary:} if some $s$--$v$ path $P$ avoids $u$, order each vertex's
    neighbours so DFS walks straight down $P$; then $v$ is discovered while $u$ is still
    undiscovered.
\end{itemize}
Equivalently: \textbf{remove $u$ and test whether $v$ is still reachable from $s$.}
The search below uses a queue, but it only tests reachability in $G\setminus\{u\}$;
any graph traversal (BFS \emph{or} DFS) works equally --- this is a reachability test,
not a DFS simulation.

\begin{minipage}{\linewidth}
\begin{Verbatim}[fontsize=\small]
MustVisitBeforeDFS(G, s, u, v):
    if v == s: return false   # s discovered first (time 0)
    if u == s: return true    # s first, so discovered before v

    visited <- array[1..n] of false
    visited[u] <- true        # "block" u: traversal never enters it

    Q <- empty queue
    visited[s] <- true
    Q.enqueue(s)
    while Q not empty:
        x <- Q.dequeue()
        for each neighbour y of x:
            if not visited[y]:
                visited[y] <- true
                Q.enqueue(y)

    # v reachable from s WITHOUT u?  No  =>  u dominates v.
    if visited[v] == false:
        return true
    else:
        return false
\end{Verbatim}
\end{minipage}

\subsection*{(b) Time complexity}

One BFS/DFS over the graph plus an $O(n)$ array initialisation: $\boxed{O(n+m)}$.

\subsection*{(c) Worked queries (source $s=A$)}

Adjacency:
\begin{verbatim}
A: B,C   B: A,D,E   C: A,G,I   D: B,F   E: B,F
F: D,E   G: C,H     H: G,I     I: H,C
\end{verbatim}

\begin{center}
\rowcolors{2}{accent!6}{white}
\begin{tabular}{l c c c p{6.4cm}}
\toprule
\rowcolor{white}
Query & $u$ & $v$ & Result & Why? ($v$ reachable from $A$ without $u$) \\
\midrule
1. $(A,C,H)$ & $C$ & $H$ & \textbf{true}  & No --- $\{G,H,I\}$ reach the rest only via $C$ \\
2. $(A,D,E)$ & $D$ & $E$ & \textbf{false} & Yes --- $A$--$B$--$E$ avoids $D$ \\
3. $(A,B,F)$ & $B$ & $F$ & \textbf{true}  & No --- $\{D,E,F\}$ reach the rest only via $B$ \\
4. $(A,H,G)$ & $H$ & $G$ & \textbf{false} & Yes --- $A$--$C$--$G$ avoids $H$ \\
\bottomrule
\end{tabular}
\end{center}

\answer{Query results: \textbf{true, false, true, false}.}

%=====================================================================
\section*{Q5. Longest path in a weighted DAG}

\subsection*{(a) Pseudocode}

Topologically sort, then DP in \textbf{reverse} topological order, where
$\mathrm{dist}[u]$ is the length of the longest path \emph{starting} at $u$, with
$\mathrm{next}[u]$ for reconstruction.

\begin{minipage}{\linewidth}
\begin{Verbatim}[fontsize=\small]
LongestPathDAG(G):
    order <- TopologicalSort(G)   # vertices in topological order
    for each v in V:
        dist[v] <- 0   # one-vertex path has length 0
        next[v] <- NULL

    for i from |order|-1 downto 0:   # reverse topological order
        u <- order[i]
        for each (v, w) in Adj[u]:
            if w + dist[v] > dist[u]:
                dist[u] <- w + dist[v]
                next[u] <- v

    best   <- the vertex maximizing dist[.]
    length <- dist[best]

    path <- empty list   # reconstruct longest path
    x <- best
    while x != NULL:
        path.append(x)
        x <- next[x]
    return (length, path)
\end{Verbatim}
\end{minipage}

\medskip
\begin{minipage}{\linewidth}
\begin{Verbatim}[fontsize=\small]
TopologicalSort(G):   # Kahn's algorithm
    compute indeg[v] for all v
    Q <- all v with indeg[v] = 0
    order <- empty list
    while Q not empty:
        u <- Q.remove()
        order.append(u)
        for each (v, w) in Adj[u]:
            indeg[v] <- indeg[v] - 1
            if indeg[v] = 0: Q.add(v)
    return order
\end{Verbatim}
\end{minipage}

\subsection*{(b) Time complexity}

Topological sort $O(n+m)$, the DP relaxes each edge once $O(n+m)$, reconstruction
$O(n)$. Total: $\boxed{O(n+m)}$.

\subsection*{(c) Worked example}

Topological order $A,B,C,D,E,F,G,H$; processing in reverse:
\begin{verbatim}
dist[H] = 0
dist[G] = 2    (G->H,2)                    next[G]=H
dist[F] = 3    (F->H,3)                    next[F]=H
dist[E] = 7    (E->H,7 > E->G:1+2)         next[E]=H
dist[D] = 4    (D->G,2: 2+2)               next[D]=G
dist[C] = 12   (C->E,5: 5+7 > C->F:4+3)    next[C]=E
dist[B] = 13   (B->E,6: 6+7 > B->D:4+4)    next[B]=E
dist[A] = 16   (A->B,3: 3+13 > A->C:2+12)  next[A]=B
\end{verbatim}

The maximum is $\mathrm{dist}[A]=16$. Reconstruct: $A\to B\to E\to H$
($3+6+7=16$).

\answer{\textbf{Longest path length $=16$}, achieved by $A\to B\to E\to H$ (edge weights $3+6+7$).}

\end{document}
